\section{Discussion}
\label{sec:discussion}

\subsection{What the tractable benchmark establishes}

The tractable benchmark is the cleanest answer to the most obvious objection to the method: does
$L_0$ only look good because it is the only method built for the largest problem? That objection is
fair, which is why Tier 1 matters. If $L_0$ performs competitively with GREG on the mixed-target
tractable benchmark and remains competitive with IPF on the count-style benchmark, then its later
advantages cannot be dismissed as a scale-only story.

This tier should also clarify what kind of advantage $L_0$ does and does not offer. If GREG
matches the tractable targets with similar accuracy but produces negative weights, that is a
different kind of limitation from simple inaccuracy. If IPF performs well on count margins but does
not extend to mixed dollar-and-count targets, that is a scope limit rather than a poor
implementation. The benchmark design is meant to keep those distinctions visible.

\subsection{What the scaling frontier establishes}

The scaling frontier shifts the question from local accuracy to method limits. A useful result in
this section is not just that one method is faster than another. It is where each method stops
being a credible tool for the target system under study. For $L_0$, the main trade-off is usually
between fit, sparsity, and runtime. For GREG, the limits are more likely to come from unstable
linear algebra, memory use, or negative weights. For IPF, the limits are more likely to come from
scope and convergence.

This frontier is important rhetorically because it distinguishes two claims that are easy to blur.
One claim is that $L_0$ can solve a hard production problem. The stronger claim is that classical
reference methods become difficult to apply before the production target volume is reached. The
scaling ladder is where that stronger claim lives.

\subsection{What the production-feasibility case establishes}

The production-feasibility tier connects the statistical benchmark back to deployment. A production
run has to do more than minimize loss on a benchmark subset. It has to produce a microdata file
that can be published, loaded, and simulated. That is where the sparsity mechanism matters most.
Without it, the cloned unit universe remains too large for many downstream uses.

This is also the tier where the paper can be most concrete about failure modes. A table that says a
method ran out of memory, returned negative weights, failed to converge, or completed in a given
runtime is more informative than a generic statement that the problem is "large." The production
case therefore should remain a feasibility table first and a performance table second.

\subsection{Trade-offs and limitations}

The benchmark does not remove all modelling judgment. Several choices remain consequential.

\subsubsection{Target-set construction}

The exact target count in a run depends on the benchmark manifest, the active target families, and
the achievable-target mask. Database-wide counts are useful for context, but the paper should base
its empirical claims on the exact exported target set for each run.

\subsubsection{Clone count and geography scope}

Clone count affects both feasibility and fidelity. More clones enlarge the feasible support for
district-level targets, but they also enlarge the optimization problem and the output dataset. The
same logic applies to geography scope. Extending the current setup to counties or other local areas
is possible in principle, but it would require new target data and a larger benchmark package.

\subsubsection{Hyperparameter sensitivity}

The $L_0$ method is not parameter free. The gate temperature, initialization, and sparsity penalty
all affect the path the optimizer takes. This is why the hyperparameter configuration belongs in an
appendix rather than disappearing into prose. The benchmark should treat those settings as part of
the method specification.

\subsubsection{Computational cost}

The pipeline is computationally heavier than a small classical calibration run. Matrix
construction, repeated policy simulation, and gradient-based optimization all cost time. That cost
is acceptable only if the pipeline solves a problem that simpler methods do not solve reliably. The
benchmark is therefore essential to justifying the added infrastructure.

\subsubsection{Comparability of IPF}

IPF is the least like-for-like baseline in the paper because its natural target representation is
different. Restricting IPF to count-style margins is the right methodological choice, but it also
means the reader should not interpret every $L_0$ versus IPF result as a full-system comparison.
The paper should say that directly.

\subsection{Generalizability}

The specific data engineering in this paper is US-specific, but the benchmark question is broader.
Many microsimulation projects need to reconcile survey microdata with administrative targets across
multiple geographic levels. Where those projects also care about keeping the resulting dataset
small enough to use, a sparse calibration method may be more useful than a classical exact-fit
procedure even if the classical procedure remains attractive on smaller subproblems.
